\section{Latar Belakang}
\label{sec:background}

\noindent Isaac Newton memandang gravitasi sebagai gaya mekanis yang menarik dua massa satu sama lain. Einstein kemudian mengubah pandangan tersebut: gravitasi bukanlah gaya, melainkan manifestasi dari kelengkungan ruang-waktu \parencite{Einstein1916}. Materi dan energi mendeformasi struktur geometri ruang-waktu, dan secara matematis kelengkungan ruang-waktu dinyatakan dalam persamaan medan Einstein.

Solusi paling sederhana dari persamaan medan Einstein dalam ruang-waktu yang vakum adalah solusi Schwarzschild, yang menggambarkan geometri di sekitar objek masif bersimetri bola \parencite{Schwarzschild1916}. Dari solusi Schwarzschild, terdapat sebuah singularitas di dalamnya, dimana kelengkungan ruang-waktu yang begitu besar menyebabkan terciptanya medan gravitasi yang sangat kuat di sekitarnya hingga partikel apapun bahkan cahaya tidak dapat lepas. Singularitas tersebut dikenal sebagai Lubang Hitam, sebuah manifestasi ruang-waktu paling ekstrem \parencite{Frolov1998}.

Karena sifat ekstremnya, lubang hitam menjadi sistem ideal untuk mempelajari lebih dalam mengenai karakteristik ruang-waktu, salah satunya adalah kestabilan lubang hitam itu sendiri. \textcite{Regge1957} mengawali teori perturbasi untuk menganalisis bagaimana respons lubang hitam Schwarzschild terhadap gangguan kecil dalam ruang vakum. Selanjutnya, \textcite{Vishveshwara1970} menemukan bahwa respons dinamis lubang hitam terhadap perturbasi dikuasai oleh mode-mode osilasi karakteristik yang disebut \textit{quasinormal modes} (QNM).

\textit{Quasinormal modes} merupakan solusi dari perturbasi yang memiliki frekuensi bersifat kompleks, dapat dituliskan sebagai $\omega = \omega_R - i\omega_I$, dengan bagian real ($\omega_R$) merepresentasikan frekuensi osilasi dan bagian imajiner ($\omega_I$) merepresentasikan laju peluruhan osilasi tersebut. Nilai dari \textit{quasinormal modes} kemudian dapat digunakan untuk menjelaskan kestabilan lubang hitam.

Dalam analisis \textit{quasinormal modes} yang telah dilakukan pada penelitian-penelitian sebelumnya, terdapat satu metodologi yang terbukti cukup akurat dan banyak digunakan untuk mendapatkan nilai frekuensi \textit{quasinormal modes}. Metode Frobenius memungkinkan perilaku solusi medan di sekitar titik-titik singular (horizon peristiwa dan tak hingga) dianalisis melalui ekspansi deret pangkat. Akan tetapi, metode Frobenius hanya dapat digunakan untuk menyatakan solusi deret di masing-masing daerah secara terpisah, sehingga diperlukan metode tambahan untuk menentukan nilai frekuensi eigen yang memenuhi kondisi batas fisik secara bersamaan di kedua daerah \parencite{Konoplya2011}. \textcite{Leaver1985} mengembangkan metode pecahan berlanjut yang memungkinkan ekstraksi frekuensi eigen \textit{quasinormal modes} dengan akurasi numerik tinggi. Kombinasi antara kedua metode tersebut banyak digunakan dalam berbagai penelitian yang telah dilakukan. Dengan demikian, penelitian ini mengaplikasikan kombinasi dari kedua metode ini untuk menentukan frekuensi eigen \textit{quasinormal modes} pada medan skalar tak bermassa maupun bermassa, khususnya pada dua latar belakang ruang-waktu: Schwarzschild dan Schwarzschild Anti-de Sitter.